\documentclass[11pt]{article}
\usepackage[margin=1in]{geometry}
\usepackage{xcolor}
\usepackage{listings}
\usepackage{enumitem}
\usepackage[colorlinks=true,urlcolor=blue,linkcolor=black]{hyperref}
\usepackage{parskip}

\definecolor{bg}{gray}{0.96}
\definecolor{cmt}{rgb}{0.30,0.45,0.30}
\lstset{
  basicstyle=\ttfamily\small,
  backgroundcolor=\color{bg},
  frame=single, framesep=6pt, rulecolor=\color{gray!40},
  breaklines=true, columns=fullflexible, keepspaces=true,
  commentstyle=\color{cmt}\itshape,
  showstringspaces=false, xleftmargin=6pt, xrightmargin=6pt,
}

\title{\vspace{-1.5em}Maritime GCS --- Terminal Launch Guide\\
\large Chattogram Port SITL + Gazebo chase-cam feed}
\author{}
\date{\today}

\begin{document}
\maketitle
\vspace{-2em}

\section{What launches}
One command brings up the whole stack. The pieces:
\begin{itemize}[leftmargin=1.4em,itemsep=2pt]
  \item \textbf{GCS web server} (\texttt{server.py}) on \texttt{http://localhost:8080} --- map, mission planner, telemetry, feed proxy.
  \item \textbf{PX4 SITL + Gazebo} (\texttt{gz\_x500}) in the \texttt{chattogram} world, home at the port.
  \item \textbf{Chase camera} (\texttt{chase\_cam.py}) --- spawns an invisible camera that follows the drone from behind/above and publishes \texttt{/chase\_cam}.
  \item \textbf{Camera bridge} (\texttt{cam\_bridge.py}) --- encodes \texttt{/chase\_cam} to MJPEG on \texttt{:8091}; the web UI's feed proxies it.
  \item \textbf{mavlink2rest} --- MAVLink $\rightarrow$ REST on \texttt{:8088} for telemetry/commands.
\end{itemize}

\section{Launch from the terminal}
\textbf{One-shot (recommended):}
\begin{lstlisting}
bash ~/maritime_gcs/start_all.sh
\end{lstlisting}
This sources the ROS 2 / Gazebo environment, kills any stale instances, then starts
the web server and \texttt{launch\_chattogram\_clean.sh} (SITL + chase cam + bridge),
all detached so they survive the terminal closing.

\textbf{Then open} \url{http://localhost:8080} and click the \texttt{Gazebo feed} button.

\medskip
\emph{Warning:} do \textbf{not} paste the launcher's contents inline. It runs
\texttt{pkill -f "cam\_bridge.py"} etc.; if those strings are in your shell's own
command line, \texttt{pkill} kills the shell. Always run it as a \emph{file}.

\section{Manual launch (if you want the pieces separately)}
\begin{lstlisting}
# 1) environment
source /opt/ros/humble/setup.bash
source ~/PX4-Autopilot/Tools/simulation/gz/setup.sh

# 2) GCS web server (:8080)
python3 ~/maritime_gcs/server.py 8080

# 3) SITL + chase cam + bridge + mavlink2rest (in another terminal)
bash ~/launch_chattogram_clean.sh
\end{lstlisting}

\section{Verify it is up ($\sim$15\,s warmup)}
\begin{lstlisting}
# web server responding
curl -s http://localhost:8080/api/kml | head -c 60

# telemetry connected (expect "connected": true, 3D fix)
curl -s http://localhost:8080/api/telemetry

# camera feed producing fresh frames
curl -s http://localhost:8091/healthz     # -> {"ok": true, "fresh": true}
\end{lstlisting}
The chase-cam feed needs $\sim$15\,s after boot (Gazebo render-engine warmup);
until then the feed shows a ``Waiting for Gazebo camera\ldots'' placeholder.

\section{Enable arming (SITL only)}
PX4 SITL blocks arming with ``No connection to the GCS''. Clear it via the pxh FIFO:
\begin{lstlisting}
for p in "COM_RC_IN_MODE 4" "CBRK_SUPPLY_CHK 894281" \
         "NAV_DLL_ACT 0" "NAV_RCL_ACT 0" "COM_RCL_EXCEPT 7"; do
  echo "param set $p" > ~/.px4_stdin; sleep 0.4
done
echo "param save" > ~/.px4_stdin
\end{lstlisting}
Full flow in the UI: \textbf{Set goal $\rightarrow$ Plan route $\rightarrow$ Upload $\rightarrow$ Arm $\rightarrow$ Start}.

\section{Stop everything}
\begin{lstlisting}
# from the UI: click "Stop SITL", or from the terminal:
pkill -f "gz sim"; pkill -f cam_bridge.py; pkill -f chase_cam.py
pkill -f mavlink2rest; pkill -f "bin/px4"
\end{lstlisting}

\section{Why the feed is a chase cam (not a screen grab)}
The old feed screen-scraped the desktop with \texttt{ffmpeg x11grab}, which returns
an all-black frame under this \textbf{Wayland} session (XWayland is rootless).
The current feed renders a real Gazebo camera on the GPU and streams it, so it works
under Wayland, Xorg, or fully headless. \texttt{chase\_cam.py} moves that camera to
track the drone every tick using the same \texttt{set\_pose} service the project's
\texttt{imu\_sim\_bridge} uses.

\section*{Ports at a glance}
\begin{tabular}{ll}
\texttt{8080} & GCS web UI + API + \texttt{/video.mjpeg} proxy \\
\texttt{8091} & chase-cam MJPEG bridge \\
\texttt{8088} & mavlink2rest REST API \\
\texttt{14550} & PX4 $\rightarrow$ mavlink2rest (UDP) \\
\end{tabular}

\end{document}
